% Weather Kernel (CS 3813)
% File: report/main.tex
% Description: Proposal/report template for the weather kernel project.
% Author: Oghenerunor Ewhro

\documentclass[11pt]{article}

\usepackage{geometry}
\geometry{margin=1in}

\usepackage{graphicx}
\usepackage{amsmath}
\usepackage{booktabs}
\usepackage{float}
\usepackage{hyperref}

\title{Project Proposal:\\
Microarchitectural Impact on Performance and Energy Efficiency\\
of a Weather Simulation Kernel}
\author{Oghenerunor Ewhro}
\date{February 2026}

\begin{document}
\maketitle

\section{Introduction}
Modern weather forecasting relies on large-scale numerical simulations that place significant
demands on processor performance, memory systems, and energy efficiency. These simulations
commonly use stencil-based finite-difference computations that repeatedly update grid points
based on neighboring values. As a result, their performance is strongly influenced by CPU
microarchitectural features such as cache hierarchies, vector execution units, and instruction-level
optimizations.

This project investigates how architectural choices and compiler optimizations affect the
performance and energy efficiency of a simplified weather simulation kernel. Rather than running
a full-scale weather model, we focus on a small representative stencil computation that captures
the core behavior of real-world simulations while remaining feasible to run on consumer hardware.

\section{Objectives}
The objectives of this project are:
\begin{itemize}
  \item To implement a simplified stencil-based weather simulation kernel in C/C++.
  \item To evaluate how compiler optimization levels affect execution time and cache behavior.
  \item To study the relationship between performance improvements and energy efficiency.
  \item To relate observed performance trends to core computer architecture concepts discussed in class.
\end{itemize}

\section{Proposed Methodology}

\subsection{Kernel Implementation}
We will implement a 2D stencil-based finite-difference kernel representative of numerical weather
prediction workloads. The kernel will update a grid over a fixed number of time steps using a local
neighborhood (e.g., a 5-point stencil). A baseline version and one optimized version will be developed
to isolate architectural effects.

\subsection{Experimental Setup}
Experiments will be conducted on consumer laptop hardware using standard, freely available tools.
We will compile and execute the kernel using different compiler optimization levels (e.g., \texttt{-O0},
\texttt{-O2}, \texttt{-O3}, and architecture-specific flags). Measurements will include execution time,
CPU utilization, cache-related metrics when available, and estimated energy consumption or power
proxies.

\subsection{Analysis}
Collected results will be analyzed to determine how changes in compiler optimizations and
microarchitectural behavior impact performance and energy efficiency. The analysis will focus on
explaining why certain configurations improve runtime and why faster execution does not always
lead to lower energy usage.

\section{Project Timeline}
The project is scheduled to begin in early March and conclude with the submission of the final
report on March 21.

\begin{center}
\begin{tabular}{ll}
\toprule
\textbf{Date Range} & \textbf{Planned Activities} \\
\midrule
Mar 1 -- Mar 7  & Baseline kernel implementation and initial measurements \\
Mar 8 -- Mar 14 & Kernel optimization and full experiment matrix \\
Mar 15 -- Mar 21 & Data analysis, report writing, and final submission \\
\bottomrule
\end{tabular}
\end{center}

\section{Division of Work and Estimated Effort}

\subsection*{Oghenerunor Ewhro}
Sole contributor responsible for implementation, experimentation, and reporting.
\begin{itemize}
  \item Baseline stencil kernel implementation
  \item Optimized kernel development
  \item Benchmark execution and data collection
  \item Data visualization and analysis
  \item Correctness validation and project documentation
\end{itemize}
Estimated effort: \textbf{48--56 hours}

\section{Expected Outcomes}
The expected outcome of this project is a clear demonstration of how CPU microarchitecture and
compiler optimizations influence the performance and energy efficiency of stencil-based scientific
computations. The final deliverables will include a written report and a presentation that connect
experimental observations to architectural principles discussed in the course.

\end{document}
